\documentclass[letterpaper]{article}
\usepackage{fontspec}
\usepackage{geometry}
\usepackage{xcolor}
\usepackage{booktabs}
\usepackage{array}
\usepackage{colortbl}
\usepackage{multicol}
\usepackage{microtype}
\usepackage{enumitem}
\usepackage{mdframed}

\geometry{letterpaper, landscape, margin=10mm}

\setmainfont{DejaVu Sans}

\definecolor{headerbg}{HTML}{2C3E50}
\definecolor{altrow}{HTML}{F0F4F8}
\definecolor{white}{HTML}{FFFFFF}
\definecolor{sectionbg}{HTML}{E8EEF4}
\definecolor{rulecol}{HTML}{C0392B}

\newcommand{\inv}[1]{\textsuperscript{#1}}
\newcommand{\rn}[2]{\textbf{#1}#2}
\newcommand{\prog}[1]{\textbf{#1}}

\setlist[itemize]{leftmargin=4mm, itemsep=1pt, topsep=2pt, parsep=0pt}
\setlist[enumerate]{leftmargin=4mm, itemsep=1pt, topsep=2pt, parsep=0pt}

\begin{document}
\pagestyle{empty}

\begin{center}
{\large\bfseries Diatonic Harp --- Progressions \& SATB Voicing Reference}
\end{center}

\vspace{2mm}

\begin{multicols}{3}

{\small\colorbox{headerbg}{\color{white}\textbf{\quad Cadences\quad}}}

\vspace{2mm}
{\small
\begin{tabular}{@{}ll@{}}
\textit{Authentic (perfect)} & \rn{V}{7}--\rn{I}{} \\[2pt]
\textit{Authentic (imperfect)} & \rn{V}{}--\rn{I}{} \\[2pt]
\textit{Half} & \rn{I}{}--\rn{V}{} or \rn{ii}{}--\rn{V}{} \\[2pt]
\textit{Plagal} & \rn{IV}{}--\rn{I}{} \\[2pt]
\textit{Deceptive} & \rn{V}{}--\rn{vi}{} \\[2pt]
\textit{Phrygian half} & \rn{iv}{}--\rn{V}{} \\
\end{tabular}
}

\vspace{4mm}
{\small\colorbox{headerbg}{\color{white}\textbf{\quad Common Progressions\quad}}}

\vspace{2mm}
{\small
\textit{Authentic close}\\
\rn{I}{}--\rn{IV}{}--\rn{V}{}--\rn{I}{}

\vspace{2mm}
\textit{Extended authentic}\\
\rn{I}{}--\rn{ii}{}--\rn{V}{}--\rn{I}{}

\vspace{2mm}
\textit{Circle of fifths}\\
\rn{I}{}--\rn{IV}{}--\rn{vii}{°}--\rn{iii}{}--\rn{vi}{}--\rn{ii}{}--\rn{V}{}--\rn{I}{}

\vspace{2mm}
\textit{Pop/hymn (I--V--vi--IV)}\\
\rn{I}{}--\rn{V}{}--\rn{vi}{}--\rn{IV}{}

\vspace{2mm}
\textit{Descending bass}\\
\rn{I}{}--\rn{I}{}\inv{1}--\rn{IV}{}--\rn{IV}{}\inv{1}--\rn{V}{}

\vspace{2mm}
\textit{Deceptive resolution}\\
\rn{I}{}--\rn{IV}{}--\rn{V}{}--\rn{vi}{}

\vspace{2mm}
\textit{ii--V--I (jazz/modal)}\\
\rn{ii}{m7}--\rn{V}{7}--\rn{I}{}Δ

\vspace{2mm}
\textit{Pachelbel / hymn}\\
\rn{I}{}--\rn{V}{}--\rn{vi}{}--\rn{iii}{}--\rn{IV}{}--\rn{I}{}--\rn{IV}{}--\rn{V}{}

\vspace{2mm}
\textit{Minor (natural)}\\
\rn{i}{}--\rn{VII}{}--\rn{VI}{}--\rn{VII}{}--\rn{i}{}

\vspace{2mm}
\textit{Minor cadential}\\
\rn{i}{}--\rn{iv}{}--\rn{V}{}--\rn{i}{}
}

\columnbreak

{\small\colorbox{headerbg}{\color{white}\textbf{\quad SATB Voice Ranges\quad}}}

\vspace{2mm}
{\small
\begin{tabular}{@{}lll@{}}
\textbf{Voice} & \textbf{Low} & \textbf{High} \\
\midrule
Soprano & C4 & G5 \\
Alto    & G3 & C5 \\
Tenor   & C3 & G4 \\
Bass    & E2 & C4 \\
\end{tabular}

\vspace{1mm}
Soprano--Alto gap: $\leq$ octave\\
Alto--Tenor gap: $\leq$ octave\\
Tenor--Bass gap: $\leq$ 12th (no limit)
}

\vspace{4mm}
{\small\colorbox{headerbg}{\color{white}\textbf{\quad Doubling Rules\quad}}}

\vspace{2mm}
{\small
\begin{itemize}
  \item Double the \textbf{root} in root-position chords
  \item Double the \textbf{bass note} in 1st inversion
  \item Avoid doubling the \textbf{leading tone} (vii / 7th scale degree)
  \item Avoid doubling the \textbf{7th} of a 7th chord
  \item In \rn{V}{7}: all four notes present --- no doubling needed
  \item \rn{vii}{°} triad: double the 3rd (not root or 5th)
  \item Diminished 5th resolves \textbf{inward}; augmented 4th resolves \textbf{outward}
\end{itemize}
}

\vspace{4mm}
{\small\colorbox{headerbg}{\color{white}\textbf{\quad Forbidden Motion\quad}}}

\vspace{2mm}
{\small
\begin{itemize}
  \item \textbf{No parallel 5ths} between any two voices
  \item \textbf{No parallel octaves} between any two voices
  \item \textbf{No direct (hidden) 5ths/octaves} in outer voices moving in similar motion
  \item \textbf{No voice crossing} (alto must stay below soprano, etc.)
  \item \textbf{No voice overlap} (alto's next note must not exceed soprano's previous note)
\end{itemize}
}

\columnbreak

{\small\colorbox{headerbg}{\color{white}\textbf{\quad Preferred Motion\quad}}}

\vspace{2mm}
{\small
\begin{itemize}
  \item \textbf{Contrary motion} between outer voices is always safe
  \item \textbf{Oblique motion} (one voice stationary) is safe
  \item \textbf{Similar motion} is acceptable if not producing parallel 5ths/8ths
  \item Prefer \textbf{stepwise motion} in all voices; avoid large leaps
  \item After a leap, resolve by \textbf{step in the opposite direction}
  \item Avoid \textbf{augmented intervals} melodically (esp. aug 2nd)
\end{itemize}
}

\vspace{4mm}
{\small\colorbox{headerbg}{\color{white}\textbf{\quad Tendency Tones\quad}}}

\vspace{2mm}
{\small
\begin{itemize}
  \item \textbf{Leading tone} (scale degree 7): always resolves \textbf{up} to tonic
  \item \textbf{Chordal 7th}: always resolves \textbf{down} by step
  \item In \rn{V}{7}--\rn{I}{}: the 7th (4̂) resolves down to 3̂; the leading tone (7̂) resolves up to 1̂
  \item \textbf{Tritone} in \rn{V}{7}: the diminished 5th (7̂--4̂) resolves inward to 1̂--3̂
\end{itemize}
}

\vspace{4mm}
{\small\colorbox{headerbg}{\color{white}\textbf{\quad Practical Tips\quad}}}

\vspace{2mm}
{\small
\begin{itemize}
  \item Keep \textbf{common tones} in the same voice when possible
  \item \rn{I}{6} (1st inversion) weakens a cadence --- use for passing chords, not finals
  \item \rn{I}{}--\rn{IV}{}: bass leaps a 4th --- move upper voices in contrary motion
  \item \rn{V}{}--\rn{vi}{} (deceptive): double the 3rd of \rn{vi}{}, not the root
  \item Avoid \textbf{unison} between adjacent voices except at cadences
  \item The \textbf{5th} is the most expendable tone --- omit before the 3rd
  \item \textbf{Open spacing} suits slow, solemn hymns; \textbf{close spacing} suits brighter passages
\end{itemize}
}

\end{multicols}

\end{document}
